\chapter{Average Case Hardness and the Overlap Gap Property}
\label{ch:ogp}

The $\NP$ completeness of the queens completion problem is a worst case statement. The experiments and exact computations of this book are on random instances, unconditioned boards or boards conditioned on queens of a solution, and they find the chain working. The gap between worst case and average case is the subject of a theory of its own, and in the last decade that theory has produced a precise geometric criterion, the overlap gap property, for when stable algorithms, local algorithms among them, must fail on random instances. This chapter presents the criterion and asks whether the queens have it.

\section{Levin's average case completeness}

Levin defined in 1986 a distributional problem as a pair of a language and a polynomial time computable distribution on instances, and showed that there are distributional problems complete for the class of all $\NP$ problems with computable distributions (extended to samplable distributions by Impagliazzo and Levin in 1990): if any one of them is solvable in polynomial time on average, all are. The theory is the formal home of Impagliazzo's Heuristica, the world in which $\NP$ problems are hard in the worst case and easy on average, and the distinction between that world and Pessiland is whether average case hard problems exist. Levin's complete problems are artificial, tiling problems with carefully chosen distributions, and the question of whether natural problems such as random satisfiability or random queens completion are hard on average is open and is attacked by different means.

\section{Random satisfiability and the clustering transition}

The natural setting is random $k$ SAT: $n$ variables, $m = \alpha n$ clauses each a random $k$ tuple of literals. For $k = 3$ the formula is believed to be satisfiable with high probability for $\alpha$ below a threshold near $4.27$ and unsatisfiable above it, by the sharp threshold theorem of Friedgut and the numerical and physical determination of the threshold by M\'ezard, Parisi and Zecchina and its rigorous confirmation for large $k$ by Ding, Sly and Sun in 2015. Below the threshold the solutions exist and the question is whether they can be found.

The statistical physics of the problem, through the cavity method, predicts a sequence of transitions in the geometry of the solution space as $\alpha$ increases: a clustering transition at which the solutions, previously a single connected set under single variable flips, shatter into exponentially many clusters separated by Hamming distance linear in $n$; a condensation transition at which a bounded number of clusters carry most of the solutions; and a freezing transition at which most variables in a typical cluster are fixed. Achlioptas and Coja Oghlan proved the clustering transition rigorously in 2008 for $k$ large enough, and for large $k$ the algorithmic barrier is observed near there: every known polynomial time algorithm finds solutions only below about the clustering density, which is near $2^k \log k / k$, a factor of order $k / \log k$ below the satisfiability threshold near $2^k \log 2$.

In the vocabulary of this book the clustering transition is the shattering of the bottom level set. Below it the solutions are predicted to form one giant component in which a local chain mixes; above it they are exponentially many components of exponentially small relative size, and no local chain crosses between them.

\section{The overlap gap property}

Gamarnik and Sudan in 2014, and Gamarnik with Jagannath, Wein and others since, turned the clustering picture into a proof technique. A problem has the overlap gap property, OGP, if there are constants $0 < \nu_1 < \nu_2 < 1$ such that, with high probability over the random instance, every pair of near optimal solutions has normalised overlap either less than $\nu_1$ or greater than $\nu_2$: there are no pairs at intermediate overlap. The property says that the solution space is clustered with a gap, in the sense that two solutions are either in the same cluster or far apart, with nothing in between.

The theorem is that OGP rules out stable algorithms. An algorithm is stable if a small change in the instance produces a small change in the output; local algorithms, low degree polynomials of the instance, approximate message passing, and Langevin dynamics run for bounded time are all stable in the appropriate sense. The argument interpolates between two instances, shows that a stable algorithm's outputs on the interpolated instances move continuously in overlap, and derives a contradiction from the gap. Gamarnik and Sudan applied it to independent sets in sparse random graphs, Gamarnik and Jagannath to the spin glass ground state, Gamarnik and Sudan and Bresler and Huang to random $k$ SAT variants, and Gamarnik and Zadik, and Wein and coauthors, to planted problems. In several cases the OGP threshold is close to the empirically observed algorithmic threshold, though not in all, and for independent sets the gap was closed only by later multi overlap arguments. For the spin glass, the absence of OGP is conjectured for the SK model and is the hypothesis under which Montanari's algorithm succeeds, and Huang and Sellke showed the matching hardness at the algorithmic threshold.

OGP is the first general tool that proves average case hardness for a natural class of algorithms without any complexity assumption, and whether the class includes the chains of this book is not known: polynomial time Markov chains are not among the algorithms for which stability has been proved. It does not prove $\Pc \neq \NP$ or anything about worst case; it proves that on random instances past the threshold, stable algorithms do not reach the optimum.

\section{Do the queens have the overlap gap property?}

Define the overlap of two queens solutions as the fraction of rows on which they agree. The question is whether, for the uniform distribution on solutions of the unconditioned problem or on completions of a randomly conditioned board, the pairwise overlaps show a gap.

For the unconditioned problem the exact computations of Chapter~\ref{ch:open} say no at small $N$: the solutions form a giant component under swaps through near solutions, and a giant component under local moves is incompatible with OGP, since consecutive states on a path within it have overlap near one and the path connects solutions of overlap near zero, filling in every intermediate value. The absence of OGP is consistent with the polynomial mixing measured in Chapter~\ref{ch:experiments2}. We conjecture that the unconditioned queens have no overlap gap for any $N$, and that this is why the Monte Carlo program works on them.

For boards conditioned on queens of a solution the computations to $N = 17$ again find a giant component and so again no gap. For adversarially conditioned boards, the reduction instances from hard formulas, the solution space is the solution space of the formula, which for a random formula above the clustering density is expected to have OGP, as proved for large $k$, and we expect, though we do not prove, that the reduction preserves the overlap structure up to the encoding. If it does, the queens completion problem has OGP on the image of hard random formulas, and the chain would fail there for the reason OGP names. This is a formal version of the expectation stated throughout the book: the hardness of the queens is the hardness of what the board encodes, and a random board encodes nothing hard.

\section{What OGP adds to the program}

The program's central conjecture is that the queens level set chain mixes torpidly on hard boards. OGP supplies a method of proof. If one can exhibit a family of boards, random in some natural sense, on which the completions have the overlap gap property, then the OGP theorem gives failure for every stable algorithm, and torpid mixing of the swap and block chains would follow if they can be shown to be stable, with no appeal to $\NP$ completeness and no complexity assumption. The reduction boards would be such a family if the overlap structure transfers through the encoding, but it would be inherited rather than intrinsic. A more satisfying result would be an intrinsic family: boards conditioned by a natural random process, such as a random fraction of fixed queens placed independently rather than taken from a solution, with a threshold in the fraction at which the completions shatter. Most of the experiments of Chapter~\ref{ch:open} fixed queens from a solution, which guarantees completions exist and biases towards connectedness. Fixing queens at random, subject only to pairwise non attack, and asking whether a completion exists and how the completions are arranged, is the natural random queens completion problem, and it has, so far as we know, not been studied before Chapter~\ref{ch:open}.

\begin{problem}
Fix $k = \kappa N$ queens on an $N \times N$ board uniformly at random subject to no two attacking, and condition on the existence of a completion. Determine the threshold $\kappa_c$ at which the completions shatter, whether the completions have the overlap gap property above it, and whether the swap chain mixes below it.
\end{problem}

This is the queens analogue of random $k$ SAT with its satisfiability and clustering thresholds, and it is, we think, the right random model in which to make the program's conjectures precise and attackable. First results at $N = 12$ are in Chapter~\ref{ch:open}: a satisfiability threshold near $\kappa = 0.4$ and a giant component that erodes to about three quarters of the completions just below it. Its solution would place the queens on the map of average case complexity that random SAT and the spin glasses already occupy, and would say in the overlap language exactly where the Monte Carlo program stops.
